\section{Price and Volume Adjustment Conventions}
\label{sec:price-volume}

Yahoo Finance's nominal \texttt{Close} history is back-adjusted for splits,
while the reported \texttt{Volume} series remains the contemporaneous share
count. Multiplying those two fields directly can understate pre-split dollar
trading value. The screen therefore identifies stock-split factors dated after
each window observation through 31 December 2025 and multiplies the vendor
close by their cumulative product. The restored close is then multiplied by
reported volume.

This is not a use of future market information to decide eligibility. The later
split factor is metadata required to undo the vendor's historical price
normalization. Future prices, returns, trading volumes, market capitalisation,
and future fundamentals are excluded from the rule. The raw selection archive
preserves the vendor close, future factor, restored close, reported volume, and
daily dollar value.

Among selected securities, \SelectedWithFutureSplitAdjustment{} require at
least one future split restoration in the 20-session window, representing
\FutureSplitEventsRestored{} distinct security-split restorations. Examples
include large post-2022 split factors; the screen restores the December 2022
price basis rather than comparing a split-normalized close with
contemporaneous volume. No selected security is assigned an unresolved
price-volume compatibility status: \UnresolvedPriceVolume{}.

Volume is interpreted as reported share volume. It is not dollar volume,
normalized volume, or a split-adjusted volume series. The research base
therefore preserves the vendor convention explicitly and avoids presenting
volume differences around splits as economic changes without an additional,
documented transformation. No ADV or relative-volume feature is constructed in
this phase.
